\documentclass[12pt]{article}

\usepackage[utf8]{inputenc}
\usepackage[margin=1in]{geometry}
\usepackage{setspace}
\doublespacing

\begin{document}

\begin{center}
{\Large\textbf{Proposal for Short Story \& Identity Memoir}}\\
\end{center}

\noindent\textbf{Student Name:} Ma Toan Bach\\
\textbf{Student Number:} 112708227\\
\textbf{Course:} LSO100KCC

\noindent\textbf{Part 1: Selected Short Story}

\noindent\textbf{Title of the Story:} ``Boys and Girls''\\
\textbf{Author:} Alice Munro

I chose Alice Munro's ``Boys and Girls'' because it shows how identity can be shaped by family expectations before a person fully understands what is happening. The narrator wants to be useful in her father's outdoor work, but she slowly learns that others expect her to become ``a girl'' in a narrower social sense.

\noindent\textbf{Part 2: Identity Focus}

My identity focus will be family expectations and personal growth. I want to explore how people learn who they are partly through the roles that family, school, and society place on them. Like Munro's narrator, I am interested in the tension between wanting to define myself and feeling pressure to become the kind of person others expect.

\noindent\textbf{Part 3: Connections Between the Story and My Memoir}

\noindent\textbf{Connection 1}

\noindent\textbf{In the story,} the narrator feels proud when she helps her father with the foxes because the outdoor work gives her a sense of importance and freedom. However, her mother and grandmother remind her that girls are supposed to behave differently, help indoors, and follow rules about how they sit, speak, and act.

\noindent\textbf{In my memoir,} I plan to connect this to a teenage memory about school and future plans, when I felt pressure to follow a path that others considered responsible. I want to reflect on the feeling of being pulled between my own goals and what other people believed was appropriate or practical for me.

\noindent\textbf{Connection 2}

\noindent\textbf{In the story,} the narrator's decision to open the gate for Flora is important because it shows a private act of resistance. Even though she cannot fully explain her action, she chooses sympathy and freedom over obedience to her father's work.

\noindent\textbf{In my memoir,} I want to connect this to the moment when I began questioning that expected path instead of simply accepting it. Like the narrator's action at the gate, the moment may seem small from the outside, but it represents the beginning of understanding my own values more clearly.

\noindent\textbf{Part 4: Preliminary Memoir Plan}

For the final memoir, I plan to discuss a memory from my teenage years when I felt pressure to follow a school or career path that others considered responsible, even though I was beginning to understand my own goals differently. The experience is meaningful because it connects to the process of learning that identity is not only something we choose individually; it is also shaped by the people around us. I will organize the memoir around this central memory, beginning with the specific situation, then reflecting on how I understood it at the time and how I understand it differently now.

\noindent\textbf{Part 5: Reflection}

At this stage, I think readers might learn that identity often develops through tension. A person may grow by recognizing expectations, questioning them, and deciding which parts to accept or resist.

\end{document}
